\textbf{COLUNI 2024 -- Questão 08}\par\medskip Leia o excerto abaixo:

\par

“Ao contar sua história, um índio, um clã, uma tribo parte do momento em que sua essência-espírito permeou a terra e relata a passagem dessa essência espírito pelos reinos vegetal, mineral e animal. Há tribos que começam a sua história desde quando o clã eram seres dos espíritos das águas. Outras trazem a sua memória animal como início da história, assim como há aquelas que iniciam a sua história a partir da árvore que foram”. (WERÁ, Kaka, 1998, p.14)

\par

Sobre o texto, é CORRETO afirmar que:\begin{enumerate}
\item constata-se a importância do registro da noção de ancestralidade, o apego às origens na definição da
      personalidade do indígena e do lugar que ocupa no mundo de hoje.
\item busca-se uma explicação para justificar o começo das narrativas contadas pelas tribos, com
      embasamentos para promover os modos de viver e de pensar dos indígenas na sua individualidade.
\item observa-se uma preocupação em dizer como se iniciam as histórias contadas por diferentes tribos, a partir
      da exemplificação dos reinos mineral, animal e vegetal.
\item apresenta-se exemplos de narrativas indígenas compostas por versões de caráteres pessoais e coletivos,
      presentes em literatura de temas indígenas.\end{enumerate}
